% Pregunta 664 (2008-julio-OEP2008, n.º 16): Mundell-Fleming, país pequeño, cambio flexible y
% movilidad perfecta del capital; expansión fiscal. Rectas a escala:
%   BP: i = i* = 2   LM: i = 0,5 + 0,6 Y   IS: i = 4 - 0,8 Y   IS': i = 5,6 - 0,8 Y
%   A = (2,5; 2)   B = IS'∩LM = (3,64; 2,69), por encima de i*: entran capitales,
%   la moneda se aprecia, caen XN y la IS' vuelve a la IS: la renta sigue en Y0.
% Autor: Víctor Gutiérrez Marcos
\documentclass[tikz,border=6pt]{standalone}
\input{tikz-tcee}
\begin{document}
\begin{tikzpicture}[tcee,x=1.15cm,y=1.15cm]
  \ejes{$Y$}{$i$}{6}{5}
  % Curvas iniciales
  \draw[curva4] (0,2) -- (5.7,2) node[etiqueta,right,text=tceedorado] {$BP$};
  \node[etiqueta,left] at (0,2) {$i^*$};
  \draw[curva1] (0.2,0.62) -- (5.6,3.86) node[etiqueta,right,text=tceemorado] {$LM$};
  \draw[curva2] (0.2,3.84) -- (4.6,0.32) node[etiqueta,right,text=tceeazul] {$IS$};
  % Lo que cambia: G sube
  \draw[curvanueva,desplazada] (2.0,4.0) -- (6.0,0.8) node[etiqueta,right,text=tceerojo] {$IS'$};
  \node[etiqueta,above,text=tceerojo] at (2.0,4.0) {$G\uparrow$};
  % Vuelta de IS' a IS por la apreciación
  \draw[flecha] (2.25,3.7) -- (0.5,3.7);
  \node[etiqueta,font=\footnotesize,align=center,text=tceerojo,anchor=north] at (1.75,3.6)
    {apreciación:\\caen $XN$};
  % Proyecciones
  \draw[guia] (2.5,2) -- (2.5,0) node[below] {$Y_0=Y_1$};
  % Puntos
  \node[punto] at (2.5,2) {};
  \node[etiqueta,above=4pt] at (2.5,2) {$A$};
  \node[punto] at (3.643,2.686) {};
  \node[etiqueta,below=4pt] at (3.643,2.686) {$B$};
  \node[etiqueta,font=\footnotesize,align=left,text=tceerojo,anchor=south west] at (3.8,2.95)
    {$B$: $i>i^*$,\\entran capitales};
\end{tikzpicture}
\end{document}
